\documentclass[instructions.tex]{subfiles}
\begin{document}
 
\chapter{Du zèle pour procurer le ciel aux âmes.}
 
\primo Si nous connoissions le prix d’une âme, que ne ferions-nous pas pour l’empêcher de périr et pour lui procurer le ciel ! C’est ce qui faisait dire à sainte Térèse : J’ai le cœur percé de douleur, quand je pense à la multitude des âmes qui périssent. Un chrétien qui aime Dieu, aime sincèrement les âmes, et voudrait les sauver toutes : il porte, pour ainsi parler, l’univers dans son cœur. Aimer les âmes pour les tirer du péché et les sauver, c’est le parfait amour du prochain. Il est bien plus avantageux d’aimer ainsi les hommes, que d’en être aimé. L’amour que les autres ont pour nous ne nous rend ni meilleurs ni plus saints; mais le zèle que nous avons pour leur salut nous rend agréables à Dieu, et nous procure le ciel.
 
On ne peut rien faire de plus avantageux pour soi-même, dit saint Eucher, que de travailler pour le bonheur d’autrui\footnote{Si amas te, proximum dilige : nihil magis commodis tuis dabis, quàm quod contuleris alienis. De verb. Sab.}. Celui, dit saint Chrysostome, qui macère son corps par les austerités, a moins de mérite que celui qui gagne une âme à Dieu; c’est quelque chose de plus grand de sauver ses frères que de faire des miracles.
 
Quand vous ne procureriez le salut que d’une seule âme, vous donneriez plus de gloire à Dieu que par vos autres bonnes œuvres, parce que vos bonnes œuvres sont imparfaites et bornées, au lieu que cette âme glorifiera Dieu parfaitement et sans fin pendant l’éternité. Par la même raison, empêcher la damnation d’une âme, c’est empêcher, en un sens, plus de mal que les hommes ne pourraient en faire pendant la durée du monde, parce qu’il n’y a point de comparaison des péchés qui se commettent en ce monde, avec les blasphèmes exécrables qu’un damné vomira contre Dieu pendant l’éternité. Il n’est donc point, dit saint Grégoire, de sacrifice plus agréable à Dieu que le zèle des âmes\footnote{Nullum Deo tale est sacrificium, quàm zelus animarum. S. Greg.}.
 
\secundo C’est une erreur de croire que les pasteurs soient les seuls obligés à procurer le salut des autres : tous peuvent et doivent y contribuer. Dieu, dit le Sage, a commandé à chacun le soin de son prochain\footnote{Mandavit unicuique de proximo suo. Eccl. 17. 12.}. Vous vous croiriez coupable, si vous laissiez périr par votre faute l’animal de votre voisin; êtes-vous moins coupable de laisser périr son âme?
 
On a du zèle pour l’entretien des hôpitaux, pour la guérison des malades : ce zèle est louable et saint; mais a-t-on le même zèle pour la guérison des âmes?
 
On s’empresse de secourir les habitans d’une ville infectée de peste, et tous les jours des âmes infectées de péchés périssent et se perdent, et nous n’y pensons pas !
 
Il n’y a personne qui ne puisse contribuer au salut des autres par quelques-uns des moyens suivans : 1. par les prières, les communions, les sacrifices. Sainte Térèse a établi son ordre particulièrement pour prier pour la conversion des pécheurs. Sainte Monique, après dix-sept ans de prières et de larmes, obtint la conversion de son mari et de son fils.
 
2. Par les pénitences. Saint François Xavier déchirait son corps par des disciplines pour obtenir la conversion d’un pécheur. Combien de personnes et de pasteurs zélés ont jeûné et pleuré plusieurs années pour obtenir le pardon à une âme !
 
3. Par l’instruction et les bons conseils. Si vous voyiez un voyageur s’égarer et prendre une route meurtrière, vous l’avertiriez. Pourquoi ne dites-vous rien à vos enfans, à vos amis, à vos compagnons, à vos compagnes, qui s’égarent dans le vice? On dirait que les hommes n’ont d’amitié et de société les uns avec les autres que pour se perdre : on n’a que trop de zèle pour le mal, pour inspirer la licence, la chicane et les querelles, et l’on ne daigne pas dire une parole pour inspirer la patience, la pudeur et la crainte de Dieu!
 
4. Par les aumônes. Quelques secours, avec un mot de salut à une pauvre famille, à quelques malades, à des prisonniers, à des chrétiens captifs chez les barbares, à qui on fait souffrir des supplices pour leur faire perdre la pudeur et la foi, arrêteraient leurs plaintes et leur désespoir, leur feraient souffrir leur infortune avec résignation. Mais, sans aller porter de secours au-delà des mers, combien parmi nous de jeunes orphelins, d’enfans abandonnés, de filles vagabondes, qui, sans retraite et sans instruction, sont dans le danger de se perdre ! Quelle charité de les retirer et de les occuper, de leur procurer l’éducation !
 
\tertio C’est ce zèle des âmes qui animait la piété d’un saint Louis, d’un saint Charles, d’une sainte Elisabeth, de tant de seigneurs et de princesses encore aujourd’hui, à la campagne et dans les villes, tant de personnes de qualité qui s’appliquent au soulagement et à la conversion de ceux que la Providence abandonne à leur charité.
 
Oh! combien d’âmes les riches pourraient-ils sauver! Les richesses leur mettent en main la clef du ciel; ils peuvent l’ouvrir pour eux-mêmes, en l’ouvrant aux autres. Que de consolations pour le moment de la mort, s’ils employaient une partie de leurs biens à faire instruire la jeunesse dans une paroisse, s’ils répandaient des livres de piété dans les familles, s’ils procuraient des ouvriers évangéliques à tant de peuples ignorans; si, par leur crédit, ils mettaient en paix des plaideurs qui se ruinent ! Oseraient-ils plaindre les dépenses qu’ils feraient pour des œuvres si saintes?
 
En vérité, l’attachement aux biens de la terre rend l’homme bien aveugle. On est avare à l’égard de Dieu, et prodigue pour le monde et pour le crime; on fait de grandes dépenses pour la vanité et le luxe, et l’on craint de faire la moindre dépense pour celui de qui on a tout reçu; rien ne coûte pour la débauche, et tout coûte quand il faut empêcher la perte d’une âme; on se fait un divertissement de hasarder de grosses sommes au jeu, et l’on épargne une obole qu’on refuse à l’instruction d’un ignorant !
 
Un repas somptueux pour des ingrats coûte plus qu’il ne faudrait pour faire instruire cent personnes, dont les prières nous ouvriraient le ciel. Un amusement de quelques heures fait souvent perdre plus d’argent qu’il n’en faudrait pour faire glorifier Dieu des années entières. Que penser d’une telle conduite dans des personnes qui se piquent d’avoir de la religion; qui, par quelques dépenses modiques, pourraient se sauver en contribuant au salut d’un millier d’âmes?
 
N’en voit-on pas même qui emploient leurs biens et leurs talens à inspirer l’erreur et l’impiété, à distribuer des livres pernicieux aux fidèles? Aurons-nous moins de zèle pour sauver nos frères, que les partisans du vice et de l’erreur n’en ont pour les pervertir?
 
On s’empresse pour obtenir des emplois dans l’Église et dans l’État, afin d’établir sa fortune; mais a-t-on le même empressement pour établir le règne de Dieu? Les peuples s’intéressent à la gloire de leur prince, on expose sa vie pour des rois mortels; mais pour la gloire du Roi du ciel, en trouve-t-on qui fassent quelques dépenses, ou qui retranchent de leurs plaisirs? Dieu veuille qu’on ne traverse pas les vues d’un pasteur qui travaille à bannir les abus; qu’on ne blâme pas son zèle, loin de le soutenir !
 
Quelle douleur pour ceux qui aiment Dieu, de voir que les intérêts de Jésus-Christ sont les seuls oubliés, et que le salut des âmes, de toutes les affaires la plus importante, est la plus négligée ! On ne devrait épargner ni dépenses, ni travaux, pour sauver les âmes; et le monde est rempli de gens qui les perdent !
 
Les apôtres s’exposaient aux tourmens pour convertir les nations; les missionnaires vont aux extrémités de la terre, au péril de leur vie, pour gagner quelques peuples à Jésus-Christ, tandis que des chrétiens riches et puissans voient d’un œil tranquille leurs frères périr; loin d’en être touchés, ils autorisent même quelquefois le libertinage et le crime.
 
Il est étrange, dit saint Bernard, qu’un animal étant tombé, on s’empresse de le relever, et que tant d’âmes étant dans le péché, personne ne leur tende la main\footnote{Cadit asinus, et invenit qui sublevet : cadit anima, et non est qui manum apponat. S. Bern.}.
 
\section{Résolutions}
 
\primo Afin d’exercer le zèle dont je suis capable, je contribuerai aux bonnes œuvres, en y consacrant ce que je pourrai de mes revenus.
 
\secundo Je répandrai quelques livres excellens, comme \textit{Pensées sur les vérités de la Religion}.
 
\tertio Je donnerai volontiers des avis à ceux qui me paraîtront en avoir besoin, mais avec douceur, saisissant les momens propices, et après avoir recommandé à Dieu mon entreprise.
 
\quarto Enfin, je prierai davantage pour la conversion des pécheurs, et j’édifierai autant que je le pourrai. 

\prayer{Mon Dieu, faites-moi la grâce de contribuer au salut de mon prochain.}
 
\end{document}
